\chapter{前言与导论}

\label{chap:00}

如果控制变量是一个向量，Boyd 的语言几乎总是够用：目标函数写在 $\mathbb R^n$ 上，约束由矩阵和超平面给出，极值、乘子和算法都能在欧氏空间里落地。但如果控制变量是函数而不是向量，或者更进一步，控制变量本身是一条轨道、一个概率测度，情况就立刻变了。此时“未知量”不再是有限个坐标，而是一个空间中的点；许多在有限维里被默认的事实，不再自动成立。

最短的桥接例子是下面这一对问题。一边是有限维二次规划
\[
\min_{x\in\mathbb R^n}\left\{\frac12 x^\top Qx+c^\top x\right\},
\]
另一边是把控制变量取为 $u\in L^2(0,T)$ 的变分问题
\[
\min_{u\in L^2(0,T)}
\left\{
\frac12\int_0^T |u(t)|^2\,dt+\Phi\!\left(\int_0^T B(t)u(t)\,dt\right)
\right\}.
\]
两者表面上都像“凸目标的极小化”，但后者已经把未知量从向量推进到了函数空间。正是在这一步，有限维直觉开始失效：闭有界集未必紧，极限不能只靠序列描述，导数和乘子必须进入对偶空间，算法也不再只是矩阵迭代，而要转化为算子分裂与邻近更新。

因此，这本书不是简单把 Boyd 的有限维结论“加一点泛函分析背景”而已，而是要把整套凸优化语言真正提升到函数、测度与随机轨道上。第 \ref{sec:0-1} 节先把“为什么必须升维”说清楚；第 \ref{sec:0-2} 节解释为什么升维以后必须同时引入拓扑、弱收敛、对偶与算子；第 \ref{sec:0-3} 节再给出不同读者的进入路径。你可以把这一章当成全书的入口，而不是作者的项目说明书。

\section{当变量不再是向量：从 Boyd 到函数空间}

\label{sec:0-1}

控制变量是函数而不是向量时，Boyd 为什么不够？更准确地说，Boyd 的有限维母语为什么不能原封不动地搬过来？原因不在于凸性失效了，而在于承载凸性的空间变了。只要未知量从 $\mathbb R^n$ 里的向量变成了一个函数、一个测度，或者一条随机轨道，问题就不再只是“坐标更多”，而是整个极限、连续性和对偶对象都换了舞台。

最短的桥接例子是把有限维二次规划和一个函数空间上的变分问题并排来看。有限维里，我们熟悉
\[
\min_{x\in\mathbb R^n}\left\{\frac12 x^\top Qx+c^\top x\right\}.
\]
如果把未知量换成控制函数 $u\in L^2(0,T)$，就会自然遇到
\[
\min_{u\in L^2(0,T)}
\left\{
\frac12\int_0^T |u(t)|^2\,dt+\Phi\!\left(\int_0^T B(t)u(t)\,dt\right)
\right\}.
\]
这两个问题共享同一层凸优化直觉：目标由一个二次项和一个复合项组成，后者都通过线性映射进入。然而它们的数学成本完全不同。前者的可行集、收敛和乘子都能在矩阵语言中处理；后者则立刻要求我们知道 $L^2(0,T)$ 的拓扑、弱收敛、对偶配对以及算子意义下的最优性条件。

\begin{example}[有限维与无穷维的最短桥]
把上面的函数空间问题投影到一个有限维基底
\[
u(t)\approx \sum_{k=1}^m \alpha_k \varphi_k(t),
\]
就会得到一个关于系数向量 $\alpha\in\mathbb R^m$ 的二次规划。反过来看，Boyd 中那些熟悉的二次规划、线性约束和投影算法，正是更一般函数空间问题在有限维离散化后的坍缩版本。把这个例子放在这里，是为了提醒读者：本书不是离开 Boyd，恰恰是在解释 Boyd 的母语来自哪里。
\end{example}

\paragraph{四个失效点}

一旦变量不再是向量，至少有四件在有限维里默认成立的事不再自动成立。

\begin{enumerate}
  \item \textbf{紧性失效}：闭有界集未必紧，极小化过程可能只得到有界而得不到强收敛极限。
  \item \textbf{极限语言升级}：弱拓扑和弱$^\ast$拓扑往往不是第一可数的，序列不再足以刻画闭包和紧性。
  \item \textbf{导数与乘子升级}：梯度不一定还是向量，最优性条件必须进入对偶空间、次微分和伴随算子。
  \item \textbf{算法接口升级}：矩阵迭代要被重新写成邻近算子、预解式和单调算子分裂。
\end{enumerate}

\begin{remark}
这四个失效点几乎就是本书后文章节安排的理由。拓扑与描述集合论解决“极限怎么说”；测度论与随机空间解决“概率对象怎么放进优化”；泛函分析与对偶理论解决“乘子和支持超平面属于哪里”；单调算子与分裂算法解决“如何把无穷维最优性系统变成可计算更新”。
\end{remark}

因此，本书的“升维”不是炫技，而是对问题本身的诚实。只要你关心的是最优控制、随机策略、图信号、核函数表示或分布式一致性，那么变量就早已不是一个向量；剩下的工作，不过是承认这一点，并为它支付必要的数学代价。第 \ref{sec:0-2} 节接下来要做的，就是把这笔代价拆开说明：为什么必须上拓扑、弱收敛、对偶与算子。Boyd 在这里仍然是重要参照，但他更像一条有限维影子，而不是终点。


\section{为什么必须上拓扑、弱收敛与对偶}

\label{sec:0-2}

上一节已经把问题摆清楚了：当变量落在函数、测度或轨道空间里时，有限维直觉并没有错，只是不再完整。接下来的问题不是“要不要学更多数学”，而是“哪一类数学分别在补哪一类缺口”。如果不把这些工具和它们所解决的问题对应起来，后文的拓扑、测度、对偶和算子章节就会显得像并列堆料。

第一个缺口是极限。极小化序列、逼近解和离散化解首先都在谈“往哪里收敛”。在有限维里，闭球紧性、序列收敛和连续映射几乎总能接上；但在无穷维里，极限常常只能在弱拓扑或弱$^\ast$拓扑下获得，于是我们必须引入拓扑语言，特别是网、闭包和下半连续性。第一个问题因此是：一个极值点究竟以什么意义存在？

第二个缺口是最优性。只要未知量不再是有限维向量，“导数 = 梯度向量”的图景就不再可靠。泛函的导数属于对偶空间，约束的法向量会变成连续线性泛函，乘子则通过伴随算子与原变量相连。因此，对偶、支撑超平面、次微分和共轭理论并不是后加的抽象层，而是无穷维最优性条件的自然语言。第二个问题于是变成：局部最优与全局对偶之间，究竟由什么对象沟通？

第三个缺口是算法。有限维中熟悉的梯度法、投影法、ADMM 和 KKT 系统，在无穷维里依然有对应物，但它们的真正母语不再是矩阵，而是邻近算子、预解式、单调算子和原--对偶分裂。算法不是在理论之后才出现的附录，它本身就在逼迫我们采用算子观点。第三个问题于是是：怎样把“求最优解”统一改写为“求某个算子的零点或不动点”？

\subsection*{记号与环境约定}

全书默认在三层环境之间切换：
\[
\text{局部凸空间} \subset \text{自反 Banach 空间} \subset \text{Hilbert 空间}
\]
这里的包含不是集合论意义，而是“结构越来越强、可用工具越来越多”的层级。讨论分离与对偶时，局部凸空间往往已经足够；讨论极小值存在和弱紧性时，通常需要自反 Banach 空间；讨论投影、梯度、邻近算子和分裂算法时，则往往要回到 Hilbert 空间。后文凡涉及对偶配对，统一写作
\[
\langle x^\ast,x\rangle.
\]
若空间本身是 Hilbert 空间，则通过 Riesz 表示可以把对偶元素重新看成向量，但这种识别必须在明确声明环境后才允许使用。

\begin{remark}
你可以把这一节看成全书的类型系统。一个命题若只在 Hilbert 空间成立，就不能被写成一般 Banach 结论；一个极小值若只在弱拓扑下存在，就不能偷换成强收敛结论；一个乘子若属于对偶空间，就不能直接画成欧氏法向量。后文所有章节的技术细化，本质上都只是把这些环境约定展开。
\end{remark}

如果把第 \ref{sec:0-1} 节的四个失效点重新对齐，那么拓扑和弱收敛负责处理极限，对偶和次微分负责处理最优性，算子和分裂负责处理算法。于是整本书的结构就不再像若干独立模块的拼装，而像对一个中心问题的三次逼近：先问极限能否存在，再问最优性如何表达，最后问如何把它算出来。第 \ref{sec:0-3} 节接下来只做一件事：告诉不同读者该沿哪条线先走。


\section{这本书该怎么读}

\label{sec:0-3}

这本书当然可以从头顺读，但未必需要从头顺读。因为它同时服务三种常见读者：想把理论基础补扎实的纯数学读者，已经懂 Boyd 母语、主要关心算法接口的读者，以及更在意建模闭环和应用落点的管理科学读者。不同入口并不意味着不同版本的书，只是意味着你先回答哪个问题。

如果你主要关心纯数学基础，那么最自然的路径是：第 \ref{chap:01} 章建立拓扑与描述集合论语言，第 \ref{chap:02} 章进入测度与随机空间，第 \ref{chap:03} 章和第 \ref{chap:04} 章建立泛函分析与分离工具，再进入第 \ref{chap:05} 章和第 \ref{chap:06} 章的次微分与对偶理论。这条线最慢，但也最稳，因为它把后文一切“算法为何有定义”“乘子为何存在”“弱极限为何足够”这些问题都提前补齐。

如果你主要关心算法，那么不必在一开始就吞下全书。更短的一条路是：先把第 \ref{chap:03} 章中的 Hilbert 空间、弱拓扑和谱论读到能用，再接第 \ref{chap:05} 章的次微分、第 \ref{chap:09} 章的单调算子与预解式，最后进入第 \ref{chap:10} 章的前向--后向分裂、Douglas--Rachford 和 ADMM。沿这条路线读时，Boyd 会一直作为有限维母语出现，但你会逐步看清：这些算法真正处理的不是矩阵，而是算子零点与原--对偶系统。

如果你主要关心管理科学、控制或应用建模，那么更合适的入口是先读随机过程与对偶接口，再回补基础。具体说，可以从第 \ref{sec:2-4} 节的随机分析入口读起，接着读第 \ref{sec:6-4} 节 Fenchel--Rockafellar 对偶、第 \ref{chap:07} 章锥优化与 KKT，再看第 \ref{chap:11} 章和第 \ref{chap:12} 章的控制、交易、图网络与推荐系统。这样你会更快看到闭环：为什么约束资格、风险对偶、分布式一致性和表示匹配都能被放回同一套凸优化语言。

无论从哪条路切入，建议都一样：把 Boyd 留在心里，但不要把它当成天花板。有限维模型是最重要的直觉来源，却不是全部世界。你越早接受“变量可以是函数、测度或轨道”，后文的拓扑、对偶和算子就越不再像门槛，而像一条自然延长线。

